\chapter{Présentation du stage}

\section{Sujet de stage}

Le sujet de stage proposé par Winamax et écrit dans la convention est le suivant:

\begin{quote}

Chez Winamax, la plateforme de jeu génère de nombreuses alertes, permettant aux personnes d'astreinte de réagir rapidement
lors d'incidents. Problématique Comment améliorer le système d'alerte, composé de plusieurs composants, afin d'offrir une solution unifiée,
moderne et modulable pour notifier efficacement les développeurs lors des incidents ? Objectifs du stage Moderniser et centraliser la stack d'alerting :
ingestion unifiée d'alertes multiples (AWS Cloudwatch, webhooks, Eventbridge…), dédoublonnage, normalisation, hiérarchisation par priorité, réduction de
la latence et du bruit, routage conditionnel, notifications multi-canaux, boucles d'escalade, historique et analytique.

\end{quote}

Compétences : RT - Mise en œuvre - Planification RT - Mise en œuvre - Développement RT - Mise en œuvre - Documentation RT - Mise en œuvre – Evaluation
RT - Exploitation - Administration RT - Exploitation - Gestion dysfonctionnement RT - Exploitation - Suivi

\section{Motivation}

Le sujet m’intéressait particulièrement, car j’ai toujours voulu travailler sur
des problématiques liées à l’observabilité.
Le fait que le projet ait pour but d’aider les développeurs au quotidien me
plaisait beaucoup.
Étant très attiré par le domaine de l’infrastructure et du DevOps, avoir l’opportunité de découvrir ce secteur au sein d’une entreprise dotée d’une infrastructure
de grande envergure représentait une véritable opportunité pour moi.

\section{Objectifs du stage}

L’objectif de mon stage était d’améliorer le système d’alerting de Winamax afin de passer
d’alarmes peu contextualisées à un modèle basé sur plusieurs niveaux de criticité, orienté vers le
service réellement impacté.